\documentclass[12pt]{article}

% ── Packages ──────────────────────────────────────────────────────────────────
\usepackage[margin=1in]{geometry}
\usepackage{graphicx}
\usepackage{float}
\usepackage{amsmath}
\usepackage{hyperref}
\usepackage{listings}
\usepackage{xcolor}
\usepackage{booktabs}
\usepackage{caption}
\usepackage{subcaption}
\usepackage{fancyhdr}
\usepackage{titlesec}
\usepackage{parskip}
\usepackage{enumitem}
\usepackage{siunitx}
\usepackage{array}
\usepackage{multirow}

% ── Code listing style ────────────────────────────────────────────────────────
\lstset{
  basicstyle=\ttfamily\small,
  keywordstyle=\color{blue!70!black},
  commentstyle=\color{gray},
  stringstyle=\color{red!70!black},
  numbers=left,
  numberstyle=\tiny\color{gray},
  frame=single,
  breaklines=true,
  captionpos=b,
  tabsize=2
}

% ── Header / Footer ───────────────────────────────────────────────────────────
\setlength{\headheight}{15pt}
\pagestyle{fancy}
\fancyhf{}
\rhead{Coherent \& Quantum Optics Lab --- Purdue ECE}
\lhead{Multi-Sensor Magnetometer PCB}
\rfoot{Page \thepage}

% ── Hyperlink colors ──────────────────────────────────────────────────────────
\hypersetup{
  colorlinks=true,
  linkcolor=blue!60!black,
  urlcolor=blue!60!black,
  citecolor=green!50!black
}

\graphicspath{{figures/}}

% ══════════════════════════════════════════════════════════════════════════════
\begin{document}

% ── Title Page ────────────────────────────────────────────────────────────────
\begin{titlepage}
  \centering
  \vspace*{2cm}
  {\LARGE\bfseries Multi-Sensor Magnetometer PCB\\[0.4em]
  \large Design, Operation, and Programming Reference\par}
  \vspace{1.5cm}
  {\large Luke Hojnicki\\
  Coherent and Quantum Optics Lab\\
  Dr.\ Elliott, Purdue University\\[0.5em]
  \today}
  \vfill
  \textit{Prepared for internal lab use. Intended audience: future lab members
  and primary user Chunya Wang.}
\end{titlepage}

% ── Table of Contents ─────────────────────────────────────────────────────────
\tableofcontents
\listoffigures
\listoftables
\newpage

% ══════════════════════════════════════════════════════════════════════════════
\section{Introduction and Purpose}

\subsection{Background}

This document describes the design, assembly, and operation of the
Magnetometer PCB developed for Dr.\ Elliott's Coherent and Quantum Optics Lab
at Purdue University. The board enables simultaneous, high-resolution
measurement of magnetic field vectors at multiple spatial positions, which is
required for Chunya Wang's experimental work characterizing magnetic field
distributions in the lab environment.

The system measures the three-dimensional magnetic field ($B_x$, $B_y$,
$B_z$) at up to eight locations using MEMSIC MMC5983MA AMR sensors arranged
across two identical PCBs connected in series via a flat flexible cable (FFC).
Each board hosts four sensors multiplexed over a single I$^2$C bus and is
controlled by a shared Arduino Nano Every microcontroller.

\subsection{Design Goals}

\begin{itemize}
  \item Simultaneous readout of 8 MMC5983MA sensors (4 per board, 2 boards)
  \item Single Arduino Nano Every host controller over a shared I$^2$C bus
  \item On-board 5\,V to 3.3\,V regulation; boards powered from the Arduino's
        USB supply
  \item Daisy-chain expansion via FFC connector (Top-on-One-Side /
        Bottom-on-the-Other cable type)
  \item Address conflict resolution between boards via hardware address
        strapping on the TCA9546A multiplexer
  \item Minimal magnetic interference from PCB-generated fields
        (see Section~\ref{sec:layout})
\end{itemize}

\subsection{Scope of This Document}

This document covers: (1) hardware overview and component descriptions,
(2) full circuit schematics, (3) PCB layout notes, (4) assembly and setup,
(5) programming the board, (6) reading and interpreting sensor data, and
(7) troubleshooting. A future student or Chunya should be able to bring the
system from unpowered to collecting data using only this document.

\newpage
% ══════════════════════════════════════════════════════════════════════════════
\section{System Architecture}

\subsection{Block Diagram}

Figure~\ref{fig:block} shows the full two-board system. The Arduino Nano Every
acts as the I$^2$C master. Both boards share the same SDA, SCL, RST, +5\,V,
and +3.3\,V lines routed through the FFC connector. Each board contains an
LM3940 voltage regulator, one TCA9546A I$^2$C multiplexer, and four
MMC5983MA sensors.

\begin{figure}[H]
  \centering
  \includegraphics[width=\textwidth]{block_diagram}
  \caption{Full system block diagram. The Arduino Nano Every drives two
  identical sensor boards (Board~0 and Board~1) via a shared FFC bus.}
  \label{fig:block}
\end{figure}

\subsection{Top-Level Schematic}

Figure~\ref{fig:sch_top} shows the top-level KiCad schematic for one board.
The FFC connector routes RST, SCL, SDA, and the four mux channel pairs
(SC0/SD0 through SC3/SD3) to the TCA9546A sub-sheet and the four MMC
sub-sheets. Two mounting holes (H1, H2) are also indicated.

\begin{figure}[H]
  \centering
  \includegraphics[width=\textwidth]{sch_toplevel}
  \caption{Top-level schematic showing FFC connector, TCA9546A, and four
  MMC5983MA instances (one board).}
  \label{fig:sch_top}
\end{figure}

\newpage
% ══════════════════════════════════════════════════════════════════════════════
\section{Component Descriptions}

\subsection{Component Summary}

\begin{table}[H]
  \centering
  \caption{Key Components (per board)}
  \begin{tabular}{llll}
    \toprule
    \textbf{Component} & \textbf{Part Number} & \textbf{Qty} &
    \textbf{Function} \\
    \midrule
    Magnetic field sensor   & MMC5983MA    & 4 & 3-axis AMR magnetometer \\
    I$^2$C multiplexer      & TCA9546AD    & 1 & 4-channel bus switch \\
    Voltage regulator       & LM3940IMPX   & 1 & 5\,V $\to$ 3.3\,V LDO \\
    FFC connector (in/out)  & \textit{[PN]} & 2 & Board-to-board bus \\
    Bulk output cap         & 33\,\textmu F & 1 & LDO output decoupling \\
    Mux bypass cap          & 10\,\textmu F & 1 & TCA9546A VCC \\
    Sensor VDD bypass cap   & 1\,\textmu F & 4 & MMC5983MA VDD \\
    Sensor CAP bypass cap   & 10\,\textmu F & 4 & MMC5983MA CAP pin \\
    I$^2$C pull-up resistors & 2.7\,k$\Omega$ & 11 & All SDA/SCL lines \\
    \bottomrule
  \end{tabular}
  \label{tab:components}
\end{table}

\subsection{MMC5983MA Magnetic Field Sensor}

The MMC5983MA (MEMSIC) is a 3-axis anisotropic magnetoresistance (AMR)
magnetic field sensor housed in a $2.0 \times 2.0 \times 1.0$\,mm LGA-16
package. It measures $B_x$, $B_y$, and $B_z$ simultaneously and communicates
over I$^2$C at a single fixed address of \texttt{0x30}. Key specifications
are listed in Table~\ref{tab:mmc_specs}.

\begin{table}[H]
  \centering
  \caption{MMC5983MA Key Specifications \cite{mmc5983ma}}
  \begin{tabular}{lll}
    \toprule
    \textbf{Parameter} & \textbf{Value} & \textbf{Notes} \\
    \midrule
    Supply voltage (VDD)       & 1.71--3.6\,V               & \\
    Full-scale range           & $\pm$8\,G                  & \\
    Resolution (18-bit mode)   & 0.0625\,mG/LSB             & 18-bit output \\
    RMS noise                  & 0.4\,mG (typ.)             & Continuous mode \\
    Output data rate           & up to \SI{1000}{\hertz}    & \\
    I$^2$C address             & \texttt{0x30} (fixed)      & \\
    Operating temperature      & $-40$ to $+85$\,\textdegree C & \\
    Package                    & LGA-16, $2.0 \times 2.0$\,mm & \\
    \bottomrule
  \end{tabular}
  \label{tab:mmc_specs}
\end{table}

\subsubsection{Pin Configuration and Connections}

In this design the MMC5983MA is used exclusively in I$^2$C mode.
Table~\ref{tab:mmc_pins} lists the relevant pin assignments; SPI-only pins
(SPL\_CS, SPL\_SDO) are left unconnected (NC).

\begin{table}[H]
  \centering
  \caption{MMC5983MA Pin Assignment (this design)}
  \begin{tabular}{lllp{6cm}}
    \toprule
    \textbf{Pin} & \textbf{Name} & \textbf{Connected to} &
    \textbf{Notes} \\
    \midrule
    1  & SCL/SPL\_SCK & SCx (from mux channel) & I$^2$C clock \\
    2  & VDD          & +3.3\,V via 1\,\textmu F cap & Supply \\
    3  & NC           & ---                    & \\
    4  & NC           & ---                    & \\
    5  & SPL\_CS      & NC                     & SPI not used \\
    6  & SPL\_SDO     & NC                     & SPI not used \\
    7--8 & NC         & ---                    & \\
    9  & GND          & GND                    & \\
    10 & CAP          & GND via 10\,\textmu F  & Internal charge pump ref \\
    11 & GND          & GND                    & \\
    12 & NC           & ---                    & \\
    13 & VDDIO        & +3.3\,V                & I/O voltage \\
    14 & NC           & ---                    & \\
    15 & INT          & NC                     & Interrupt not used \\
    16 & SDA/SPL\_SDI & SDx (from mux channel) & I$^2$C data \\
    \bottomrule
  \end{tabular}
  \label{tab:mmc_pins}
\end{table}

\subsubsection{SET/RESET Coil}

The MMC5983MA contains an on-chip SET/RESET coil that drives a large current
pulse through the sensing element to re-magnetize the Permalloy film in a
known direction. This is necessary because:

\begin{itemize}
  \item \textbf{SET} magnetizes the film in the ``positive'' direction,
        establishing a known reference state and removing any negative offset
        from strong external field exposure.
  \item \textbf{RESET} magnetizes the film in the opposite direction.
        Taking the difference of a SET and RESET measurement,
        $B = (B_\text{SET} - B_\text{RESET})/2$, cancels the
        temperature-dependent bridge offset, substantially improving accuracy.
\end{itemize}

A SET or RESET pulse is issued by writing the corresponding bit in the
\texttt{CONTROL\_0} register (\texttt{0x09}). The firmware must wait at
least \SI{1}{\micro\second} after the pulse before issuing a measurement
command. See Section~\ref{sec:setresetproc} for the recommended procedure.

\subsubsection{Fixed I\textsuperscript{2}C Address and the Multiplexer Requirement}

Because all MMC5983MA devices share the fixed I$^2$C address \texttt{0x30},
placing more than one sensor on a single bus causes address collisions.
The TCA9546A multiplexer resolves this by isolating each sensor on its own
switched bus segment, presenting only one sensor to the master at a time.

\subsection{TCA9546A I\textsuperscript{2}C Multiplexer}
\label{subsec:tca}

The TCA9546AD (Texas Instruments) is a low-voltage 4-channel I$^2$C bus
switch in a VSSOP-10 package. It sits between the Arduino Nano Every (I$^2$C
master) and the four MMC5983MA sensors, enabling the master to address each
sensor individually by enabling only one downstream channel at a time.

\subsubsection{I\textsuperscript{2}C Address Strapping}

The TCA9546A I$^2$C address is determined by the logic levels on address
pins A0--A2 (pins 1, 2, 13 in the schematic). The base address is
\texttt{0b1110\_xxx} where \texttt{xxx = A2A1A0}, giving the range
\texttt{0x70}--\texttt{0x77}.

On this board, A2 and A1 are tied to GND. A0 is the address differentiation
pin used to distinguish Board~0 from Board~1:

\begin{table}[H]
  \centering
  \caption{TCA9546A Address Strapping per Board}
  \begin{tabular}{lllll}
    \toprule
    \textbf{Board} & \textbf{A2} & \textbf{A1} & \textbf{A0} &
    \textbf{I$^2$C Address} \\
    \midrule
    Board 0 & GND & GND & GND & \texttt{0x70} \\
    Board 1 & GND & GND & VCC & \texttt{0x71} \\
    \bottomrule
  \end{tabular}
  \label{tab:tca_addr}
\end{table}

\textbf{Important:} Both boards ship with A0 tied to GND (\texttt{0x70}).
When connecting two boards in series, the A0 pin on Board~1's TCA9546A
must be moved from GND to VCC to assign it address \texttt{0x71}. This
is a solder bridge or jumper change on the board. See
Section~\ref{sec:twoboard} for the two-board setup procedure.

\subsubsection{Channel Select Register}

The TCA9546A is controlled by writing a single byte to its I$^2$C address.
Bits [3:0] enable channels 0--3 respectively; only one bit should be set
at a time. Writing \texttt{0x00} disconnects all channels (recommended
after each read to prevent bus contention).

\begin{table}[H]
  \centering
  \caption{TCA9546A Channel Select Byte}
  \begin{tabular}{lll}
    \toprule
    \textbf{Channel} & \textbf{Byte (hex)} & \textbf{Sensor} \\
    \midrule
    0 & \texttt{0x01} & U4 / MMC\_0 \\
    1 & \texttt{0x02} & U3 / MMC\_1 \\
    2 & \texttt{0x04} & U5 / MMC\_2 \\
    3 & \texttt{0x08} & U6 / MMC\_3 \\
    None & \texttt{0x00} & All disconnected \\
    \bottomrule
  \end{tabular}
  \label{tab:tca_select}
\end{table}

\subsubsection{Pull-up Resistors}

Each I$^2$C segment has dedicated pull-up resistors to +3.3\,V. The values
used are $2.7\,\text{k}\Omega$ throughout (R1--R11), consistent with
\SI{400}{\kilo\hertz} Fast-Mode operation at the expected bus capacitances.
Specifically:

\begin{itemize}
  \item R1, R2, R3: upstream SDA, SCL, and one additional line
  \item R4/R5, R6/R7, R8/R9, R10/R11: downstream SDA/SCL pairs for
        channels 0--3 (SD0/SC0 through SD3/SC3)
\end{itemize}

\subsubsection{Bypass Capacitor}

A 10\,\textmu F capacitor (C1) is placed on the TCA9546A VCC pin for
bulk decoupling.

\subsubsection{RESET Pin}

The TCA9546A \texttt{/RESET} pin is connected to the RST net from the FFC
connector, which is driven by the Arduino Nano Every. Asserting this pin
low disconnects all channels and resets the channel select register to
\texttt{0x00}. This allows the firmware to perform a hardware bus reset
without power-cycling the entire board.

\subsection{LM3940IMPX Voltage Regulator}

The LM3940IMPX (U2) is a low-dropout (LDO) linear regulator that converts
the 5\,V supply from the Arduino Nano Every USB rail to the 3.3\,V required
by the TCA9546A and MMC5983MA sensors. A 33\,\textmu F output capacitor
(C2) provides output stabilization and transient suppression as specified
in the LM3940 datasheet. Both GND pins of the regulator are connected to
the board ground plane. \textbf{Do not power the sensor VDD from the
Arduino's 5\,V pin directly.}

\subsection{FFC Connector and Daisy-Chain Bus}

Each board has two FFC connectors (J4 PIN\_IN, J1 PIN\_OUT) for daisy-chain
expansion. The 16-pin connector carries the following signals:

\begin{table}[H]
  \centering
  \caption{FFC Connector Pinout}
  \begin{tabular}{lll}
    \toprule
    \textbf{Pin(s)} & \textbf{Signal} & \textbf{Notes} \\
    \midrule
    1, 2, 15, 16 & GND       & Ground \\
    3            & +5V       & Arduino USB rail \\
    7            & SCL       & I$^2$C clock (shared bus) \\
    10           & SDA       & I$^2$C data (shared bus) \\
    12           & RST       & Active-low mux reset \\
    14           & +3.3V     & Regulated output from LM3940 \\
    4--6, 8--9, 11, 13 & NC  & Not connected \\
    \bottomrule
  \end{tabular}
  \label{tab:ffc_pins}
\end{table}

\textbf{Cable orientation note:} The FFC cable used in this system is of the
``Top-on-One-Side, Bottom-on-the-Other'' (Type B) variety. This means the
contacts on one end face up and the contacts on the other end face down.
Ensure the cable is inserted into each connector in the correct orientation
to avoid a pin-reversal short. When connecting Board~0 PIN\_OUT to
Board~1 PIN\_IN, verify that power and signal nets align before applying
power.

\subsection{Arduino Nano Every}

The Arduino Nano Every serves as the I$^2$C master and USB-serial interface.
It is based on the ATmega4809 microcontroller, operates at \SI{20}{\mega\hertz},
and provides 3.3\,V-compatible I/O.

\begin{table}[H]
  \centering
  \caption{Arduino Nano Every Key Specifications}
  \begin{tabular}{ll}
    \toprule
    \textbf{Parameter} & \textbf{Value} \\
    \midrule
    Microcontroller    & ATmega4809 \\
    CPU speed          & \SI{20}{\mega\hertz} \\
    Flash / SRAM       & 48\,kB / 6\,kB \\
    I$^2$C pins        & A4 (SDA), A5 (SCL) \\
    I$^2$C speed       & Standard (100\,kHz) / Fast (400\,kHz) \\
    USB interface      & USB CDC via SAMD11 co-processor \\
    \bottomrule
  \end{tabular}
  \label{tab:nano_specs}
\end{table}

The Nano Every's A4/A5 pins connect to the shared SDA/SCL bus. A digital
output pin connects to the RST net for hardware mux resets. The board is
powered from USB; the on-board 3.3\,V regulator output is also forwarded to
the sensor boards via the FFC connector.

Upload instructions: select \textbf{Arduino Nano Every} under
\textit{Tools $\rightarrow$ Board $\rightarrow$ Arduino megaAVR Boards},
set \textit{Registers Emulation} to \textbf{None (ATMEGA4809)}, select
the correct COM port, and click Upload.

\newpage
% ══════════════════════════════════════════════════════════════════════════════
\section{Circuit Schematics}

\subsection{Power Distribution}
\label{sec:power_sch}

Figure~\ref{fig:sch_power} shows the power and connector schematic. The
LM3940IMPX (U2) receives +5\,V from the FFC connector (J4 PIN\_IN, pin 3)
and outputs +3.3\,V (VCC net) stabilized by a 33\,\textmu F capacitor (C2).
The regulated +3.3\,V is also forwarded to the PIN\_OUT connector (J1) so
that it is available to downstream boards in the daisy chain. SCL, SDA,
RST, and GND are passed through both connectors unchanged.

\begin{figure}[H]
  \centering
  \includegraphics[width=\textwidth]{sch_power}
  \caption{Power distribution and FFC connector schematic. U2 (LM3940IMPX)
  converts +5\,V to +3.3\,V (VCC). J4 is the input FFC connector;
  J1 passes signals to the next board in the chain.}
  \label{fig:sch_power}
\end{figure}

\subsection{I\textsuperscript{2}C Multiplexer (TCA9546A)}

Figure~\ref{fig:sch_mux} shows the TCA9546A multiplexer schematic. The
upstream SDA and SCL lines from the FFC connector enter pins 15 and 14
of U1 respectively, with $2.7\,\text{k}\Omega$ pull-ups (R1--R3) to VCC.
Each of the four downstream channel pairs (SD0/SC0 through SD3/SC3) has
its own $2.7\,\text{k}\Omega$ pull-up pair to +3.3\,V (R4--R11). Address
pins A2 and A1 (pins 13 and 2) are tied to GND; A0 (pin 1) determines
the board address (GND = \texttt{0x70}, VCC = \texttt{0x71}). The RESET
pin (pin 3) is driven by the RST net from the FFC connector.

\begin{figure}[H]
  \centering
  \includegraphics[width=0.95\textwidth]{sch_mux}
  \caption{TCA9546AD I$^2$C multiplexer schematic. All pull-ups are
  $2.7\,\text{k}\Omega$ to +3.3\,V. A0 straps select the board I$^2$C
  address (see Table~\ref{tab:tca_addr}).}
  \label{fig:sch_mux}
\end{figure}

\subsection{MMC5983MA Sensor Instances}

All four sensors per board are identical circuits; they differ only in which
mux channel they connect to (SC0/SD0 through SC3/SD3). Each sensor has:
a 1\,\textmu F bypass capacitor on VDD (pins C5, C3, C7, C9 for sensors
0--3) and a 10\,\textmu F capacitor on the CAP pin (C6, C4, C8, C10) to
GND, as required by the MMC5983MA datasheet. SPL\_CS and SPL\_SDO are left
unconnected to select I$^2$C mode. The INT pin is not connected as interrupt-
driven acquisition is not used in this implementation.

\begin{figure}[H]
  \centering
  \begin{subfigure}[b]{0.48\textwidth}
    \includegraphics[width=\textwidth]{sch_mmc0}
    \caption{U4 / MMC\_0 (channel 0, SC0/SD0)}
  \end{subfigure}
  \hfill
  \begin{subfigure}[b]{0.48\textwidth}
    \includegraphics[width=\textwidth]{sch_mmc1}
    \caption{U3 / MMC\_1 (channel 1, SC1/SD1)}
  \end{subfigure}
  \\[1em]
  \begin{subfigure}[b]{0.48\textwidth}
    \includegraphics[width=\textwidth]{sch_mmc2}
    \caption{U5 / MMC\_2 (channel 2, SC2/SD2)}
  \end{subfigure}
  \hfill
  \begin{subfigure}[b]{0.48\textwidth}
    \includegraphics[width=\textwidth]{sch_mmc3}
    \caption{U6 / MMC\_3 (channel 3, SC3/SD3)}
  \end{subfigure}
  \caption{MMC5983MA sensor schematics. All four instances are identical;
  each is isolated on its own TCA9546A downstream channel.}
  \label{fig:sch_sensors}
\end{figure}

\newpage
% ══════════════════════════════════════════════════════════════════════════════
\section{PCB Layout}
\label{sec:layout}

\subsection{Board Overview}

The PCB measures \textbf{50.05\,mm $\times$ 120.05\,mm} and uses a
\textbf{2-layer stackup} (F\_Cu / B\_Cu) on a standard 1.6\,mm FR4 substrate.
It was designed in KiCad~10.0.1. The board outline has a distinctive
hourglass profile: a wider octagonal control section at the top, a narrow
routing neck in the middle, and a wider octagonal sensor section at the
bottom with a large circular cutout at its center.

\begin{table}[H]
  \centering
  \caption{PCB Fabrication Specifications}
  \begin{tabular}{ll}
    \toprule
    \textbf{Parameter} & \textbf{Value} \\
    \midrule
    Board dimensions    & 50.05\,mm $\times$ 120.05\,mm \\
    Layer count         & 2 (F\_Cu top, B\_Cu bottom) \\
    Board thickness     & 1.6\,mm \\
    Min trace width     & 0.3\,mm \\
    Pad-to-pad clearance & 0.2\,mm \\
    Pad-to-track clearance & 0.2\,mm \\
    Surface finish      & \textit{[HASL / ENIG / specify]} \\
    \bottomrule
  \end{tabular}
  \label{tab:fab_specs}
\end{table}

\subsection{Top and Bottom Copper Layers}

\begin{figure}[H]
  \centering
  \begin{subfigure}[b]{0.45\textwidth}
    \includegraphics[width=\textwidth]{pcb_top}
    \caption{Top copper (F\_Cu, red) with silkscreen.}
  \end{subfigure}
  \hfill
  \begin{subfigure}[b]{0.45\textwidth}
    \includegraphics[width=\textwidth]{pcb_bottom}
    \caption{Bottom copper (B\_Cu, blue).}
  \end{subfigure}
  \caption{PCB layout rendered from Gerber files. Top layer carries nearly
  all routing; the bottom layer provides a partial ground plane and return
  paths.}
  \label{fig:pcb_layers}
\end{figure}

\subsection{Board Shape and Functional Regions}

The board is divided into three functional regions:

\begin{enumerate}
  \item \textbf{Control section (top octagon):} Houses the two FFC
        connectors, the LM3940IMPX regulator (U2) with its 33\,\textmu F
        output cap (C2), the TCA9546A multiplexer (U1) with its 10\,\textmu F
        bypass cap (C1), and all 11 pull-up resistors (R1--R11). Two
        mounting holes are located at the top corners of this section.

  \item \textbf{Routing neck (center):} A narrow bridge carrying the eight
        I$^2$C channel traces (SD0--SD3, SC0--SC3) from the multiplexer
        to the sensor section. Keeping these traces confined to the neck
        minimizes the copper area in the magnetically sensitive sensor
        region.

  \item \textbf{Sensor section (bottom octagon):} Contains the four
        MMC5983MA sensors (U3--U6) and their associated bypass capacitors
        (C3--C10), arranged around a large circular cutout
        (approximately 30\,mm radius / 60\,mm diameter). The cutout allows
        the board to be mounted concentrically around a cylindrical
        experimental object such as a glass cell or optical tube.
\end{enumerate}

\subsection{Sensor Placement Around the Circular Cutout}

The four sensors are positioned at the cardinal points of the circular
cutout to provide spatially distributed field measurements around the
mounted object:

\begin{table}[H]
  \centering
  \caption{Sensor Positions Around Circular Cutout}
  \begin{tabular}{llll}
    \toprule
    \textbf{Ref.} & \textbf{Channel} & \textbf{Clock position} &
    \textbf{Notes} \\
    \midrule
    U3 & 1 (SC1/SD1) & 12 o'clock (top)    & \\
    U4 & 0 (SC0/SD0) & 9 o'clock (left)    & \\
    U6 & 3 (SC3/SD3) & 3 o'clock (right)   & \\
    U5 & 2 (SC2/SD2) & 6 o'clock (bottom)  & \\
    \bottomrule
  \end{tabular}
  \label{tab:sensor_positions}
\end{table}

Each sensor is oriented with its package aligned to the board axes. If
absolute field direction relative to the mounted object is required,
the orientation of each sensor's $X$, $Y$, $Z$ axes relative to the
board frame should be recorded when installing the board in the
experimental apparatus.

\subsection{Magnetic Noise Minimization}

The following layout decisions were made to minimize PCB-generated magnetic
interference at the sensor locations:

\begin{itemize}
  \item \textbf{Functional isolation:} All active components that carry
        significant current (regulator, mux, connectors) are confined to
        the control section, physically separated from the sensors by the
        narrow neck. This puts the largest current loops as far as possible
        from the sensor LGAs.
  \item \textbf{Bypass cap proximity:} Each sensor's 1\,\textmu F VDD cap
        and 10\,\textmu F CAP pin cap are placed immediately adjacent to the
        sensor pads, minimizing the loop area of high-frequency switching
        currents.
  \item \textbf{No ferromagnetics near sensors:} All passive components in
        the sensor section are ceramic capacitors; no inductors or
        ferrite-core components are placed in the bottom octagon.
  \item \textbf{Trace routing in neck:} The I$^2$C channel traces are
        bundled tightly through the neck and routed as differential-pair-like
        groups (SDA/SCL per channel adjacent) to reduce loop area and
        radiated field.
  \item \textbf{Silkscreen note:} The board carries a ``BOILER UP!!''
        silkscreen label on the control section --- this is cosmetic only
        and has no electrical significance.
\end{itemize}

\newpage
% ══════════════════════════════════════════════════════════════════════════════
\section{Assembly and Setup}

\subsection{Bill of Materials (BOM)}

\begin{table}[H]
  \centering
  \caption{Bill of Materials (one board)}
  \begin{tabular}{lllll}
    \toprule
    \textbf{Ref.} & \textbf{Value / PN} & \textbf{Qty} &
    \textbf{Function} & \textbf{Supplier} \\
    \midrule
    U1          & TCA9546AD      & 1 & I$^2$C mux       & \textit{[TI / Digikey]} \\
    U2          & LM3940IMPX     & 1 & 5V$\to$3.3V LDO  & \textit{[TI / Digikey]} \\
    U3--U6      & MMC5983MA      & 4 & 3-axis sensor    & \textit{[MEMSIC / Digikey]} \\
    J1, J4      & FFC 16-pin     & 2 & Board connector  & \textit{[PN]} \\
    C1          & 10\,\textmu F  & 1 & TCA VCC bypass   & \textit{[PN]} \\
    C2          & 33\,\textmu F  & 1 & LDO output cap   & \textit{[PN]} \\
    C3, C5, C7, C9 & 1\,\textmu F & 4 & Sensor VDD bypass & \textit{[PN]} \\
    C4, C6, C8, C10 & 10\,\textmu F & 4 & Sensor CAP pin & \textit{[PN]} \\
    R1--R11     & 2.7\,k$\Omega$ & 11 & I$^2$C pull-ups & \textit{[PN]} \\
    H1, H2      & Mounting hole  & 2  & Mechanical      & --- \\
    \bottomrule
  \end{tabular}
  \label{tab:bom}
\end{table}

\subsection{Soldering Notes}

\begin{itemize}
  \item The MMC5983MA and TCA9546AD are LGA/VSSOP packages that require
        reflow soldering. Hand-soldering is possible for the VSSOP but
        not recommended for the LGA sensors.
  \item Apply flux generously under LGA pads before placing components.
  \item After reflow, inspect all sensor and mux connections with a
        microscope or loupe. Tombstoned or bridged pads will cause
        I$^2$C failures.
  \item The LM3940IMPX is a through-hole or SMD package (verify footprint);
        solder last to avoid heat damage to nearby sensors.
\end{itemize}

\subsection{Initial Power-On Checklist}

\begin{enumerate}
  \item \textbf{Before connecting the Arduino:} Measure the +3.3\,V output
        of U2 with a DMM. It should read $3.3\,\text{V} \pm 0.1\,\text{V}$.
        If it reads 0\,V or 5\,V, check U2 orientation and C2 polarity.
  \item Verify that SDA and SCL have \SI{3.3}{\volt} pull-ups by measuring
        them at rest (should read $\approx3.3\,\text{V}$ with no activity).
  \item Connect the Arduino and upload the I$^2$C scanner sketch
        (Section~\ref{sec:i2c_scan}).
  \item Confirm that \texttt{0x70} (Board 0 mux) appears in the scan.
        If running two boards, also confirm \texttt{0x71} (Board 1 mux).
  \item Select each mux channel in turn and verify that \texttt{0x30}
        (MMC5983MA) appears on each channel. Eight total hits expected
        for the two-board system.
  \item Run the test firmware and confirm plausible field readings
        ($\pm$ambient Earth field, roughly 0.2--0.6\,G depending on
        orientation).
\end{enumerate}

\newpage
% ══════════════════════════════════════════════════════════════════════════════
\section{Two-Board Setup}
\label{sec:twoboard}

\subsection{Overview}

Two identical boards are connected in series via an FFC cable to provide
eight sensor channels total. The Arduino Nano Every drives both boards over
the same SDA, SCL, RST, and power lines. The two TCA9546A multiplexers are
differentiated by I$^2$C address:

\begin{itemize}
  \item \textbf{Board 0:} Mux address \texttt{0x70} (A0 = GND,
        factory default)
  \item \textbf{Board 1:} Mux address \texttt{0x71} (A0 = VCC,
        requires hardware modification)
\end{itemize}

\subsection{Hardware Modification for Board 1}

On Board 1's TCA9546A (U1), the A0 pin must be changed from GND to VCC:

\begin{enumerate}
  \item Identify the A0 pad on U1 (pin 1 of TCA9546AD; see
        Figure~\ref{fig:sch_mux}).
  \item Remove or cut the GND connection to A0.
  \item Bridge A0 to the VCC net (3.3\,V). A short wire or solder
        bridge to the adjacent VCC pad is sufficient.
  \item After the modification, confirm with an ohmmeter that A0 is
        connected to VCC and isolated from GND.
\end{enumerate}

\subsection{FFC Cable Connection}

\begin{enumerate}
  \item Use a Type-B FFC cable (Top-on-One-Side, Bottom-on-the-Other).
        This type is required because the PIN\_OUT connector on one board
        faces opposite to the PIN\_IN connector on the next.
  \item Insert one end of the cable into Board 0's PIN\_OUT connector (J1).
        The contacts should face the correct direction for the ZIF/LIF latch
        on that connector.
  \item Insert the other end into Board 1's PIN\_IN connector (J4).
  \item \textbf{Before powering on:} double-check that the +5\,V, +3.3\,V,
        GND, SDA, and SCL pins map correctly across the connector. A pin
        reversal will short power to GND.
\end{enumerate}

\subsection{Sensor Numbering Convention}

With both boards connected, the eight sensors are addressed as follows:

\begin{table}[H]
  \centering
  \caption{Full 8-Sensor Addressing Map}
  \begin{tabular}{lllll}
    \toprule
    \textbf{Logical ID} & \textbf{Board} & \textbf{Mux addr} &
    \textbf{Channel} & \textbf{Sensor I$^2$C addr} \\
    \midrule
    Sensor 0 & Board 0 & \texttt{0x70} & 0 (\texttt{0x01}) & \texttt{0x30} \\
    Sensor 1 & Board 0 & \texttt{0x70} & 1 (\texttt{0x02}) & \texttt{0x30} \\
    Sensor 2 & Board 0 & \texttt{0x70} & 2 (\texttt{0x04}) & \texttt{0x30} \\
    Sensor 3 & Board 0 & \texttt{0x70} & 3 (\texttt{0x08}) & \texttt{0x30} \\
    Sensor 4 & Board 1 & \texttt{0x71} & 0 (\texttt{0x01}) & \texttt{0x30} \\
    Sensor 5 & Board 1 & \texttt{0x71} & 1 (\texttt{0x02}) & \texttt{0x30} \\
    Sensor 6 & Board 1 & \texttt{0x71} & 2 (\texttt{0x04}) & \texttt{0x30} \\
    Sensor 7 & Board 1 & \texttt{0x71} & 3 (\texttt{0x08}) & \texttt{0x30} \\
    \bottomrule
  \end{tabular}
  \label{tab:sensor_map}
\end{table}

\newpage
% ══════════════════════════════════════════════════════════════════════════════
\section{Programming the Board}

\subsection{Development Environment}

\begin{itemize}
  \item Arduino IDE 2.x (or 1.8.x)
  \item Board package: \textit{Arduino megaAVR Boards} (install via
        Boards Manager, search ``megaAVR'')
  \item Required libraries: \texttt{Wire} (built-in),
        \textit{[MMC5983MA library --- add name and source here]}
\end{itemize}

\subsection{I\textsuperscript{2}C Scanner}
\label{sec:i2c_scan}

Use this sketch to verify hardware before running the full firmware:

\begin{lstlisting}[language=C++, caption={I$^2$C scanner sketch}]
#include <Wire.h>

void setup() {
  Wire.begin();
  Wire.setClock(400000);  // 400 kHz Fast Mode
  Serial.begin(115200);
  Serial.println("I2C Scanner");
}

void loop() {
  for (uint8_t addr = 1; addr < 127; addr++) {
    Wire.beginTransmission(addr);
    if (Wire.endTransmission() == 0) {
      Serial.print("Found device at 0x");
      Serial.println(addr, HEX);
    }
  }
  Serial.println("-- scan complete --");
  delay(3000);
}
\end{lstlisting}

Expected output (two boards):
\begin{lstlisting}
Found device at 0x70   <- Board 0 mux
Found device at 0x71   <- Board 1 mux
\end{lstlisting}

After confirming both muxes, scan inside each channel:

\begin{lstlisting}[language=C++, caption={Per-channel I$^2$C scan}]
void scanChannel(uint8_t muxAddr, uint8_t channel) {
  // Select channel
  Wire.beginTransmission(muxAddr);
  Wire.write(1 << channel);
  Wire.endTransmission();

  // Scan
  Wire.beginTransmission(0x30);
  if (Wire.endTransmission() == 0)
    Serial.println("  MMC5983MA found at 0x30");
  else
    Serial.println("  No sensor!");

  // Deselect
  Wire.beginTransmission(muxAddr);
  Wire.write(0x00);
  Wire.endTransmission();
}
\end{lstlisting}

\subsection{Core Firmware Structure}

\begin{lstlisting}[language=C++, caption={Firmware skeleton --- main loop}]
#include <Wire.h>
// #include <MMC5983MA.h>  // add your library here

#define MUX0  0x70
#define MUX1  0x71
#define MMC   0x30
#define RST_PIN 4   // Arduino pin connected to RST net

void selectSensor(uint8_t muxAddr, uint8_t channel) {
  Wire.beginTransmission(muxAddr);
  Wire.write(1 << channel);
  Wire.endTransmission();
}

void deselectAll(uint8_t muxAddr) {
  Wire.beginTransmission(muxAddr);
  Wire.write(0x00);
  Wire.endTransmission();
}

void resetMuxes() {
  digitalWrite(RST_PIN, LOW);
  delayMicroseconds(10);
  digitalWrite(RST_PIN, HIGH);
}

void readSensor(uint8_t muxAddr, uint8_t channel,
                float &Bx, float &By, float &Bz) {
  selectSensor(muxAddr, channel);

  // TODO: issue SET pulse, trigger measurement, read 18-bit registers
  // Convert raw counts to Gauss:
  //   B_G = (raw - 2^17) / (2^18 / 16.0)

  deselectAll(muxAddr);
}

void setup() {
  Wire.begin();
  Wire.setClock(400000);
  Serial.begin(115200);
  pinMode(RST_PIN, OUTPUT);
  digitalWrite(RST_PIN, HIGH);
  resetMuxes();
}

void loop() {
  float Bx, By, Bz;
  // Print header
  Serial.println("sensorID,Bx_G,By_G,Bz_G");

  // Board 0 (mux 0x70), sensors 0-3
  for (uint8_t ch = 0; ch < 4; ch++) {
    readSensor(MUX0, ch, Bx, By, Bz);
    Serial.print(ch); Serial.print(",");
    Serial.print(Bx,4); Serial.print(",");
    Serial.print(By,4); Serial.print(",");
    Serial.println(Bz,4);
  }

  // Board 1 (mux 0x71), sensors 4-7
  for (uint8_t ch = 0; ch < 4; ch++) {
    readSensor(MUX1, ch, Bx, By, Bz);
    Serial.print(ch+4); Serial.print(",");
    Serial.print(Bx,4); Serial.print(",");
    Serial.print(By,4); Serial.print(",");
    Serial.println(Bz,4);
  }

  delay(100);  // 10 Hz acquisition
}
\end{lstlisting}

\newpage
% ══════════════════════════════════════════════════════════════════════════════
\section{Reading and Interpreting Data}

\subsection{Output Format}

The firmware streams CSV over USB-serial at 115200 baud. Each line
represents one sensor reading:

\begin{lstlisting}
sensorID,Bx_G,By_G,Bz_G
0,0.2341,0.0123,-0.4412
1,0.2339,0.0120,-0.4410
...
7,0.2344,0.0125,-0.4415
\end{lstlisting}

Open the Arduino Serial Monitor or connect a Python script
(see Section~\ref{sec:dataworkflow}).

\subsection{Unit Conversion}

The MMC5983MA outputs 18-bit unsigned integers. The zero-field output
code is $2^{17} = 131072$. The full-scale range is $\pm8$\,G across
$2^{18} = 262144$ counts, giving a sensitivity of
$16\,\text{G} / 262144 = 0.0625$\,mG/LSB. The conversion is:

\begin{equation}
  B\,[\text{G}] = \frac{\text{raw} - 2^{17}}{2^{18}/16}
               = \frac{\text{raw} - 131072}{16384}
  \label{eq:conversion}
\end{equation}

where raw is the 18-bit output value (combined from the six measurement
registers \texttt{0x00}--\texttt{0x05}, plus two LSBs in \texttt{0x06}).

\subsection{SET/RESET Procedure}
\label{sec:setresetproc}

Issue a SET/RESET at the beginning of each acquisition session and
periodically during long runs (recommended every 5--10 minutes, or if
readings drift unexpectedly):

\begin{enumerate}
  \item Write \texttt{0x08} to \texttt{CONTROL\_0} (\texttt{0x09}) to
        issue a \textbf{SET} pulse. Wait $\geq1\,\mu$s.
  \item Take measurement $B_\text{SET}$.
  \item Write \texttt{0x10} to \texttt{CONTROL\_0} to issue a
        \textbf{RESET} pulse. Wait $\geq1\,\mu$s.
  \item Take measurement $B_\text{RESET}$.
  \item Compute $B = (B_\text{SET} - B_\text{RESET})/2$ to cancel the
        DC offset.
\end{enumerate}

For continuous monitoring where offset cancellation is less critical,
a single SET at startup followed by triggered measurements is acceptable.

\subsection{Data Collection Workflow (for Chunya)}
\label{sec:dataworkflow}

\begin{enumerate}
  \item Connect the Arduino Nano Every to the computer via USB.
  \item Open a terminal or run the Python logger script:
\begin{lstlisting}[language=Python, caption={Minimal Python logger}]
import serial, csv, time

port = "COM3"          # Windows: COMx; Linux/Mac: /dev/ttyACM0
baud = 115200

with serial.Serial(port, baud, timeout=1) as ser, \
     open("data.csv", "w", newline="") as f:
    writer = csv.writer(f)
    writer.writerow(["time_s","id","Bx","By","Bz"])
    t0 = time.time()
    while True:
        line = ser.readline().decode().strip()
        if line and not line.startswith("sensor"):
            parts = line.split(",")
            writer.writerow([round(time.time()-t0,3)] + parts)
\end{lstlisting}
  \item Allow the system to warm up for 30 seconds before recording data.
  \item Issue a SET/RESET at the start of each run (see above or
        press the reset button to re-run \texttt{setup()}).
  \item Stop logging with Ctrl-C; data is saved to \texttt{data.csv}.
  \item Verify data with \texttt{pandas}: confirm all 8 sensor IDs
        appear and field magnitudes are in the range 0.2--0.6\,G
        (ambient Earth field).
\end{enumerate}

\newpage
% ══════════════════════════════════════════════════════════════════════════════
\section{Troubleshooting}

\begin{table}[H]
  \centering
  \caption{Common Issues and Resolutions}
  \begin{tabular}{p{3.8cm} p{4.2cm} p{5cm}}
    \toprule
    \textbf{Symptom} & \textbf{Likely Cause} & \textbf{Resolution} \\
    \midrule
    I$^2$C scanner shows nothing
      & Pull-ups missing; solder joint on mux
      & Verify 2.7\,k$\Omega$ on SDA/SCL; reflow U1 \\
    \addlinespace
    Only one mux found (e.g., only \texttt{0x70})
      & Board 1 A0 not moved to VCC
      & Rework A0 solder bridge on Board 1 TCA9546A (see Section~\ref{sec:twoboard}) \\
    \addlinespace
    Mux found but sensor \texttt{0x30} missing on a channel
      & Bad reflow on that MMC5983MA; wrong channel byte
      & Re-solder suspect sensor; check channel select logic \\
    \addlinespace
    All-zeros or maxed-out readings
      & SET/RESET not issued; sensor not in I$^2$C mode (SPL\_CS floating)
      & Issue SET before measurement; verify SPL\_CS is NC not floating \\
    \addlinespace
    Noisy / drifting readings
      & Insufficient SET/RESET; ferromagnetic object nearby
      & Repeat SET/RESET; remove tools or cables from sensor vicinity \\
    \addlinespace
    FFC cable connects but no comms
      & Cable orientation wrong (not Type-B)
      & Flip cable end-for-end; re-check contact orientation \\
    \addlinespace
    Upload fails (Nano Every)
      & Wrong board/Registers Emulation setting
      & Set Registers Emulation to ``None (ATMEGA4809)''; double-tap reset \\
    \addlinespace
    LM3940 output $\neq$ 3.3\,V
      & C2 wrong polarity; U2 misorientated
      & Verify C2 ($+$ to OUT pin); check U2 pin 1 marker \\
    \bottomrule
  \end{tabular}
  \label{tab:troubleshoot}
\end{table}

\subsection{Using an Oscilloscope to Debug I\textsuperscript{2}C}

Probe SDA on pin 15 and SCL on pin 14 of U1. A healthy transaction at
400\,kHz should show:
\begin{itemize}
  \item SCL: clean square wave, 400\,kHz period, logic-high $\approx3.3\,\text{V}$.
  \item SDA: data transitions only when SCL is low; stable during SCL high.
  \item ACK: SDA pulled low by the device on the 9th clock pulse.
\end{itemize}

If SDA does not return high between transactions, suspect a stuck device.
Assert the RST pin low for $>10\,\mu$s to reset the mux and clear any
stuck state.

\subsection{Bus Recovery Procedure}

If the I$^2$C bus hangs (SDA stuck low):
\begin{enumerate}
  \item Assert RST\_PIN low from the Arduino for \SI{100}{\micro\second}.
  \item Release RST\_PIN high.
  \item Call \texttt{Wire.end()} then \texttt{Wire.begin()} in firmware to
        reinitialize the TWI peripheral.
  \item Re-run the I$^2$C scanner to confirm bus recovery.
\end{enumerate}

\newpage
% ══════════════════════════════════════════════════════════════════════════════
\section{Known Limitations and Future Work}

\begin{itemize}
  \item \textbf{No on-chip temperature compensation:} The MMC5983MA offset
        drifts with temperature. For precision measurements over long runs,
        periodically repeat the SET/RESET subtraction procedure.
  \item \textbf{No interrupt-driven acquisition:} The INT pin is unused;
        the firmware polls the measurement-done bit. A future revision
        could wire INT to an Arduino input to reduce latency and CPU load.
  \item \textbf{Expansion limited to 8 addresses (0x70--0x77):} A maximum
        of 8 boards (32 sensors) can share a single I$^2$C bus using the
        TCA9546A address range. Beyond this, a second hardware I$^2$C
        bus or a different mux topology is required.
  \item \textbf{FFC cable length sensitivity:} Long FFC cables increase
        bus capacitance. If communication errors occur with cables
        $>$30\,cm, reduce \texttt{Wire.setClock()} to 100\,kHz.
  \item \textbf{SPI mode not implemented:} The MMC5983MA supports SPI for
        faster data rates ($>$1\,MHz). If the 400\,kHz I$^2$C throughput
        proves insufficient for future experiments, a board redesign
        using SPI would allow much higher sample rates.
\end{itemize}

\newpage
% ══════════════════════════════════════════════════════════════════════════════
\section*{Acknowledgments}
\addcontentsline{toc}{section}{Acknowledgments}

Thank you to Dr.\ Elliott for guidance throughout this project, and to
Chunya Wang for testing the two-board configuration and identifying
the I$^2$C address conflict (and its solution) during system bring-up.

\newpage
% ══════════════════════════════════════════════════════════════════════════════
\begin{thebibliography}{9}

\bibitem{mmc5983ma}
  MEMSIC, \textit{MMC5983MA 3-Axis Magnetic Sensor with I$^2$C/SPI Interface
  Datasheet}, Rev.\ A.

\bibitem{tca9546a}
  Texas Instruments, \textit{TCA9546A Low-Voltage 4-Channel I$^2$C Switch
  with Reset Datasheet}, SCPS204.

\bibitem{lm3940}
  Texas Instruments, \textit{LM3940 1-A Low-Dropout Regulator for 5-V to
  3.3-V Conversion Datasheet}, SNOS386.

\bibitem{nanoevery}
  Arduino, \textit{Arduino Nano Every Datasheet and Pinout Reference},
  \url{https://docs.arduino.cc/hardware/nano-every}.

\bibitem{kicad}
  KiCad EDA, \url{https://www.kicad.org}.

\end{thebibliography}

% ══════════════════════════════════════════════════════════════════════════════
\appendix

\section{I\textsuperscript{2}C Address Map}

\begin{table}[H]
  \centering
  \caption{Complete I$^2$C Address Map (two-board system)}
  \begin{tabular}{llll}
    \toprule
    \textbf{Device} & \textbf{Address} & \textbf{Board} &
    \textbf{Notes} \\
    \midrule
    TCA9546A (U1) & \texttt{0x70} & Board 0 & A0=GND \\
    TCA9546A (U1) & \texttt{0x71} & Board 1 & A0=VCC \\
    MMC5983MA (U3--U6) & \texttt{0x30} & Both & Isolated per mux channel \\
    Arduino Nano Every & --- & --- & Master only, no slave addr \\
    \bottomrule
  \end{tabular}
  \label{tab:i2c_map}
\end{table}

\section{Register Quick Reference}

\begin{table}[H]
  \centering
  \caption{MMC5983MA Key Registers}
  \begin{tabular}{lllp{6cm}}
    \toprule
    \textbf{Reg (hex)} & \textbf{Name} & \textbf{R/W} & \textbf{Description} \\
    \midrule
    0x00 & Xout0  & R & X-axis output [17:10] \\
    0x01 & Xout1  & R & X-axis output [9:2] \\
    0x02 & Yout0  & R & Y-axis output [17:10] \\
    0x03 & Yout1  & R & Y-axis output [9:2] \\
    0x04 & Zout0  & R & Z-axis output [17:10] \\
    0x05 & Zout1  & R & Z-axis output [9:2] \\
    0x06 & XYZout2 & R & Bits [1:0] of X, Y, Z (18-bit) \\
    0x08 & Status & R & Bit 0: Meas\_M\_Done \\
    0x09 & Control\_0 & W & Bit 0: TM\_M (trigger); Bit 3: SET; Bit 4: RESET \\
    0x0A & Control\_1 & W & BW[1:0]: bandwidth / ODR select \\
    0x2F & Product ID & R & Should read \texttt{0x30} \\
    \bottomrule
  \end{tabular}
  \label{tab:registers}
\end{table}

\section{Gerber File Manifest}

All Gerber files were generated from KiCad~10.0.1 on 2026-04-27.
The job file is \texttt{MAG\_PCB\_FINAL-job.gbrjob}.

\begin{table}[H]
  \centering
  \caption{Gerber File Manifest}
  \begin{tabular}{lll}
    \toprule
    \textbf{Filename} & \textbf{Layer} & \textbf{Function} \\
    \midrule
    \texttt{MAG PCB FINAL-F\_Cu.gbr}         & F\_Cu (L1)       & Top copper, all signal routing \\
    \texttt{MAG PCB FINAL-B\_Cu.gbr}         & B\_Cu (L2)       & Bottom copper, ground/return paths \\
    \texttt{MAG PCB FINAL-F\_Paste.gbr}      & F\_Paste         & Top solder paste (stencil) \\
    \texttt{MAG PCB FINAL-B\_Paste.gbr}      & B\_Paste         & Bottom solder paste \\
    \texttt{MAG PCB FINAL-F\_Silkscreen.gbr} & F\_Silkscreen    & Top component labels and outline \\
    \texttt{MAG PCB FINAL-B\_Silkscreen.gbr} & B\_Silkscreen    & Bottom labels \\
    \texttt{MAG PCB FINAL-F\_Mask.gbr}       & F\_Mask          & Top solder mask (openings) \\
    \texttt{MAG PCB FINAL-B\_Mask.gbr}       & B\_Mask          & Bottom solder mask \\
    \texttt{MAG PCB FINAL-Edge\_Cuts.gbr}    & Edge\_Cuts       & Board outline / routed profile \\
    \texttt{MAG PCB FINAL-job.gbrjob}        & ---              & Gerber job file (fab instructions) \\
    \bottomrule
  \end{tabular}
  \label{tab:gerbers}
\end{table}

Submit all ten files to the PCB fabricator. The job file encodes the
board size (50.05\,mm $\times$ 120.05\,mm), layer count (2),
thickness (1.6\,mm), and design rules automatically.
\end{document}
